\begin{frame}{두 가지 차이: 고정이냐 재해석이냐}
  \begin{columns}[T]
    \begin{column}{0.5\textwidth}
      \kb{차이 1: 문맥 고정 vs 재해석}
      \begin{itemize}
        \item word2vec: 벡터는 \textbf{문맥과 무관하게 고정} ---
          같은 단어는 \textbf{항상} 같은 벡터.
        \item Transformer: 토큰 임베딩은 뒤의 \textbf{어텐션 층들을
          거치며 문맥에 따라 계속 재해석} --- ``은행''의 초기
          임베딩은 하나지만, 문장에 따라 최종 표현은 \textbf{다름}.
      \end{itemize}
    \end{column}
    \begin{column}{0.5\textwidth}
      \kb{차이 2: 산출물이냐 조각이냐}
      \begin{itemize}
        \item word2vec: \textbf{벡터 자체가 최종 산출물}.
        \item Transformer: 토큰 임베딩은 \textbf{훨씬 큰 모델의 한
          조각}일 뿐.
      \end{itemize}
    \end{column}
  \end{columns}
  \vskip1em
  \kq{12장에서 ``가정''했던 것이, 14.2의 word2vec과 같은 정신으로
  ``학습된 임베딩''이었다.}
\end{frame}
